% Chapter 8: Exact Learning and Query Models
% Centerpiece: Angluin's L* algorithm and the passive-active gap.
% DFAs are NOT PAC-learnable (Kearns-Valiant 1994) but ARE exactly
% learnable with MQ+EQ (Angluin 1987). This separation is the content.

\chapter{Exact Learning and Query Models}
\label{ch:exact}

The previous chapters study learners that receive data passively: drawn from a distribution (PAC), or presented by an adversary (online), or enumerated by nature (Gold). This chapter studies what happens when the learner can \emph{ask questions}.

The exact learning model, introduced by Angluin~\cite{Angluin1987}, equips the learner with two oracles: a membership query oracle and an equivalence query oracle. The learner's goal is not approximate generalization but \emph{exact identification}: the game ends when the learner proposes a hypothesis that the equivalence oracle accepts. The central result of this chapter is that DFAs, computationally intractable under passive data~\cite{KearnsValiant1994}, become efficiently learnable with query access. This gap between passive and active learning is one of the sharpest structural results in the field.

% ================================================================
% SECTION 1: THE EXACT LEARNING FRAMEWORK
% ================================================================

\section{The Exact Learning Framework}
\label{sec:exact-framework}

The query model was introduced in \Cref{sec:queries}, where membership and equivalence oracles were defined (\Cref{def:mq,def:eq}). We recall the framework here with the precision needed for algorithmic analysis.

\begin{defn}[Minimally Adequate Teacher]
\label{def:mat}
A \emph{minimally adequate teacher} (MAT) for a target concept $c$ over domain $X$ provides two oracles:
\begin{itemize}[leftmargin=*, itemsep=2pt]
\item $\mathrm{MQ}(x)$: returns $c(x)$ for any $x \in X$ chosen by the learner.\formal{FLT_Proofs.Data.MembershipOracle}
\item $\mathrm{EQ}(h)$: if $h = c$, returns \textsc{yes}. Otherwise, returns a counterexample $z \in X$ with $h(z) \neq c(z)$.\formal{FLT_Proofs.Data.EquivalenceOracle}
\end{itemize}
The teacher is ``minimally adequate'' in the sense that it provides exactly these two capabilities and nothing more: no distributional information, no structural hints, no preference among counterexamples.\formal{FLT_Proofs.Learner.Active.MinimallyAdequateTeacher}
\end{defn}

\begin{defn}[Exact Learning]
\label{def:exact-learning-formal}
A concept class $\mathcal{C}$ is \emph{exactly learnable in polynomial time} if there exists an algorithm $A$ such that for every target $c \in \mathcal{C}$, algorithm $A$ interacts with a MAT for $c$ and eventually receives \textsc{yes} from $\mathrm{EQ}$, using a number of queries and computation steps polynomial in:
\begin{enumerate}[leftmargin=*, nosep]
\item the representation size $n$ of the target $c$, and
\item the length $m$ of the longest counterexample returned by $\mathrm{EQ}$.
\end{enumerate}\formal{FLT_Proofs.Criterion.PAC.ExactLearnable}
\end{defn}

\begin{remark}[The role of counterexample length]
\label{rmk:counterexample-length}
The dependence on $m$ is necessary because the equivalence oracle's counterexamples are adversarial: the oracle may return arbitrarily long counterexamples. The algorithm must process them, so the complexity must account for their length.
\end{remark}

% ================================================================
% SECTION 2: THE L* ALGORITHM
% ================================================================

\section{Angluin's $L^*$ Algorithm}
\label{sec:lstar}

The $L^*$ algorithm learns the minimal DFA for an unknown regular language $L \subseteq \Sigma^*$ using membership and equivalence queries. It is the canonical example of exact learning and the centerpiece of this chapter.

\subsection{Observation Tables}

The algorithm maintains an \emph{observation table} $(S, E, T)$, where:
\begin{itemize}[leftmargin=*, itemsep=2pt]
\item $S \subseteq \Sigma^*$ is a finite, prefix-closed set of \emph{access strings} (candidate state representatives).
\item $E \subseteq \Sigma^*$ is a finite, suffix-closed set of \emph{distinguishing extensions}.
\item $T : (S \cup S \cdot \Sigma) \times E \to \{0,1\}$ is the table, where $T(s, e) = \mathrm{MQ}(s \cdot e)$.
\end{itemize}
For each string $s \in S \cup S \cdot \Sigma$, define the \emph{row} $\mathrm{row}(s) = (T(s, e_1), T(s, e_2), \ldots, T(s, e_{|E|}))$.

\begin{defn}[Closed and Consistent Table]
\label{def:closed-consistent}
The observation table $(S, E, T)$ is:
\begin{itemize}[leftmargin=*, itemsep=2pt]
\item \emph{Closed} if for every $s \in S$ and $a \in \Sigma$, there exists $s' \in S$ with $\mathrm{row}(s \cdot a) = \mathrm{row}(s')$.
\item \emph{Consistent} if for all $s_1, s_2 \in S$ with $\mathrm{row}(s_1) = \mathrm{row}(s_2)$, we have $\mathrm{row}(s_1 \cdot a) = \mathrm{row}(s_2 \cdot a)$ for all $a \in \Sigma$.
\end{itemize}
\end{defn}

Closedness ensures that every one-step extension of a state representative is equivalent to some existing representative, so the transition function of the conjectured DFA is well-defined. Consistency ensures that equivalent representatives behave identically under extension, so the transition function is well-defined as a function of equivalence classes, not of particular strings.

\subsection{The Algorithm}

\begin{figure}[ht]
\centering
\begin{tikzpicture}[
    >={Stealth[length=2.2mm]},
    pole/.style={draw, very thick, rounded corners, align=center, font=\small\bfseries,
                 minimum width=2.3cm, minimum height=1.3cm},
    lstep/.style={draw, semithick, rounded corners, align=center, font=\scriptsize,
                 fill=sepgreen!10, draw=sepgreen!55!black, minimum width=3.0cm,
                 minimum height=0.9cm, inner sep=3pt},
    cyc/.style={-{Stealth[length=2mm]}, semithick, defblue!60!black},
    q/.style={-{Stealth[length=2mm]}, semithick, thmred!70!black},
    lit/.style={draw, densely dashed, thmred!70!black, rounded corners, fill=thmred!4,
                align=left, font=\scriptsize, text=black!72, inner sep=4pt}
]
% --- agents (poles) ---
\node[pole, fill=defblue!10, draw=defblue!60!black] (learner) at (-0.3,0)
   {Learner\\[-1pt]{\scriptsize\mdseries table $(S,E,T)$}};
\node[pole, fill=thmred!10, draw=thmred!60!black, minimum height=1.6cm] (mat) at (11,0)
   {MAT\\[-1pt]{\scriptsize\mdseries knows $A^{\ast}$}\\[-1pt]{\tiny\mdseries\leanname{MinimallyAdequateTeacher}}};
% --- the L* state-discovery loop (verified data-type nodes) ---
\node[lstep] (n1) at (3.3,2.1) {1.~Fill table via $\mathrm{MQ}$\\{\tiny\leanname{MembershipOracle}}};
\node[lstep] (n2) at (3.3,0) {2.~Close \& make consistent\\{\tiny\cref{def:closed-consistent}}};
\node[lstep] (n3) at (3.3,-2.1) {3.~Conjecture DFA $M$\\{\tiny\leanname{DFA'}}};
\node[lstep] (n4) at (7.7,-2.1) {4.~$\mathrm{EQ}(M)$\\{\tiny\leanname{EquivalenceOracle}}};
\node[lstep] (n5) at (7.7,2.1) {5.~counterexample $z$\\{\tiny\leanname{Counterexample}}};
% --- cycle edges ---
\draw[cyc] (n1) -- (n2);  \draw[cyc] (n2) -- (n3);  \draw[cyc] (n3) -- (n4);
\draw[cyc] (n4) -- (n5);  \draw[cyc] (n5) -- (n1);
\draw[cyc, densely dashed] (n2.west) to[out=158,in=202]
   node[font=\tiny, left, align=center]{not closed/\\consistent} (n1.west);
% --- query interaction with the MAT ---
\draw[q] (learner.north) to[out=65,in=115] node[font=\scriptsize, above]{membership queries $\mathrm{MQ}$} (mat.north);
\draw[q] (n4.east) to[out=0,in=205] node[font=\tiny, above right, pos=0.6]{$\mathrm{EQ}$} (mat.south west);
\draw[q, densely dashed] (mat.west) to[out=155,in=0] node[font=\tiny, above, pos=0.45]{$z$ / \textsc{yes}} (n5.east);
% --- the criterion (Will) ---
\node[font=\tiny, text=defblue!55!black, align=center, anchor=north] at (3.3,-2.85)
   {halt iff $\mathrm{EQ}=$\textsc{yes} (\leanname{ExactLearnable})};
% --- literature: correctness + termination (Angluin 1987, not kernel) ---
\node[lit, text width=3.5cm, anchor=north] at (7.9,-2.9)
   {\textbf{Correctness \& the $n-1$ bound:} converges to the minimal DFA in $\le n-1$ equivalence queries ($n=$ Myhill--Nerode index). Angluin 1987 (\cref{thm:lstar-complexity}); open in the kernel (\cref{op:lstar-formal}).};
\end{tikzpicture}
\caption[The $L^*$ exact-learning loop and its Myhill--Nerode reading]{The $L^*$ protocol as the
exact-learning-via-a-MAT loop. A learner and a minimally adequate teacher cooperate to identify the
minimal DFA for an unknown regular language. The observation table $(S,E,T)$ is a finite approximation of
the Myhill--Nerode right-congruence of the target: its distinct rows are the equivalence classes
discovered so far. Membership queries fill the table cell by cell; an equivalence query returns a
counterexample that forces a new state class, splitting the current approximation. Closedness and
consistency are exactly the conditions under which the table induces a well-defined DFA, and termination
in at most $n-1$ equivalence queries is the index of the congruence ($n$ states in the minimal DFA). The
data-types are verified scaffolding: the membership and equivalence oracles
(\leanname{MembershipOracle}, \leanname{EquivalenceOracle}), the teacher
(\leanname{MinimallyAdequateTeacher}), the algorithm object (\leanname{LStar}), the conjecture type
(\leanname{DFA'}), and the success criterion (\leanname{ExactLearnable}). The correctness of the loop and
the $n-1$ bound are Angluin's theorem (\cite{Angluin1987}), cited and not yet kernel-proved
(\cref{op:lstar-formal}); they are drawn dashed.}
\label{fig:lstar-protocol}
\end{figure}

The $L^*$ algorithm proceeds as follows.

\begin{enumerate}[leftmargin=*, itemsep=4pt]
\item \textbf{Initialize.} Set $S = E = \{\varepsilon\}$ (the empty string). Fill the table $T$ using membership queries.

\item \textbf{Close and make consistent.} Repeat until the table is both closed and consistent:
\begin{itemize}[nosep]
\item If the table is not closed (there exists $s \in S$, $a \in \Sigma$ with $\mathrm{row}(s \cdot a) \neq \mathrm{row}(s')$ for any $s' \in S$), add $s \cdot a$ to $S$ and fill new table entries via $\mathrm{MQ}$.
\item If the table is not consistent (there exist $s_1, s_2 \in S$ with $\mathrm{row}(s_1) = \mathrm{row}(s_2)$ but $\mathrm{row}(s_1 \cdot a) \neq \mathrm{row}(s_2 \cdot a)$ for some $a \in \Sigma$), then there exists $e \in E$ witnessing the difference. Add $a \cdot e$ to $E$ and fill new entries via $\mathrm{MQ}$.
\end{itemize}

\item \textbf{Conjecture.} When the table is closed and consistent, construct a DFA $M$:
\begin{itemize}[nosep]
\item States: the distinct rows $\mathrm{row}(s)$ for $s \in S$.
\item Initial state: $\mathrm{row}(\varepsilon)$.
\item Transition: $\delta(\mathrm{row}(s), a) = \mathrm{row}(s \cdot a)$ (well-defined by closedness and consistency).
\item Accepting states: those $\mathrm{row}(s)$ with $T(s, \varepsilon) = 1$.
\end{itemize}

\item \textbf{Query.} Submit $\mathrm{EQ}(M)$.
\begin{itemize}[nosep]
\item If the teacher returns \textsc{yes}, halt and output $M$.
\item If the teacher returns a counterexample $z$, add $z$ and all its prefixes to $S$, fill the table, and return to step~2.
\end{itemize}
\end{enumerate}

\begin{thm}[$L^*$ correctness and complexity {\cite{Angluin1987}}]
\label{thm:lstar-complexity}
Let $A^*$ be the minimal DFA for the target language $L$, with $n$ states over alphabet $\Sigma$. The $L^*$ algorithm terminates and outputs the minimal DFA for $L$, using:
\begin{itemize}[nosep]
\item at most $n - 1$ equivalence queries, and
\item at most $O(n^2 |\Sigma| \cdot m)$ membership queries,
\end{itemize}
where $m$ is the length of the longest counterexample received.
\end{thm}

\begin{proof}[Proof sketch]
Each counterexample introduces at least one new distinct row in the observation table, because the counterexample witnesses a string that the current conjecture misclassifies. Since the minimal DFA has $n$ states and the table's distinct rows correspond to Myhill--Nerode equivalence classes, the number of distinct rows is bounded by $n$. The table starts with one row and gains at least one per counterexample, so at most $n - 1$ equivalence queries occur before the table has $n$ distinct rows, at which point the conjecture DFA is isomorphic to $A^*$.

The membership query bound follows from the table dimensions: $|S| \leq n$, $|S \cdot \Sigma| \leq n|\Sigma|$, and $|E| \leq m \cdot n$ (each counterexample adds at most $m$ suffixes). The total number of cells is $O(n|\Sigma|) \times O(nm) = O(n^2 |\Sigma| m)$, each requiring one membership query.

The full correctness proof (that a closed, consistent table yields a DFA consistent with all membership queries, and that the Myhill--Nerode theorem guarantees uniqueness) is given in Angluin~\cite{Angluin1987}.
\end{proof}

\begin{remark}[Why the table works]
\label{rmk:table-intuition}
The observation table is a finite approximation to the Myhill--Nerode equivalence relation of the target language. Two strings $s_1, s_2$ are Myhill--Nerode equivalent if $s_1 \cdot w \in L \iff s_2 \cdot w \in L$ for all $w \in \Sigma^*$. The table approximates this by testing only the extensions in $E$. Closedness ensures the approximation is fine enough to define transitions; consistency ensures it is coherent. Each counterexample reveals that the current approximation conflates two genuinely distinct equivalence classes, forcing a refinement.
\end{remark}

% ================================================================
% SECTION 3: THE PASSIVE-ACTIVE GAP
% ================================================================

\section{The Passive--Active Gap}
\label{sec:passive-active-gap}

The structural significance of $L^*$ lies not in the algorithm itself but in what it implies when combined with a hardness result.

\begin{thm}[DFAs are not PAC-learnable {\cite{KearnsValiant1994}}]
\label{thm:dfa-pac-hard}
Under standard cryptographic assumptions (specifically, the hardness of inverting RSA or factoring Blum integers), the class of DFAs with $n$ states is not efficiently PAC-learnable from random examples alone. No polynomial-time algorithm can PAC-learn DFAs, even improperly.
\end{thm}

The proof reduces the problem of breaking a cryptographic primitive to the problem of PAC-learning DFAs: if a polynomial-time PAC learner existed, it could be used to invert the cryptographic function, contradicting the hardness assumption. The reduction is non-trivial and we defer to Kearns and Valiant~\cite{KearnsValiant1994} for the full argument.

\begin{separation}
\textbf{Passive $\not\Leftrightarrow$ Active.} Witness: the class of DFAs.
\begin{itemize}[leftmargin=*, nosep]
\item \textbf{Active learning succeeds.} $L^*$ exactly identifies any $n$-state DFA using $O(n^2 |\Sigma| m)$ membership queries and $n - 1$ equivalence queries (\Cref{thm:lstar-complexity}).
\item \textbf{Passive learning fails.} Under cryptographic assumptions, no polynomial-time algorithm PAC-learns DFAs from random examples (\Cref{thm:dfa-pac-hard}).
\end{itemize}
The gap is not quantitative (more data needed) but \emph{qualitative} (no amount of passive data suffices for efficient learning, while a polynomial number of queries does). The mechanism is clear: membership queries let the learner probe the target's transition structure at chosen points, and equivalence queries provide directed feedback at the learner's current frontier of error. Random examples provide neither capability.
\end{separation}

This separation witnesses a fundamental edge in the concept graph: $\nde{pac\_learning} \not\Rightarrow \nde{exact\_learning}$ and $\nde{exact\_learning} \not\Rightarrow \nde{pac\_learning}$. Neither paradigm subsumes the other. The reverse direction (classes that are PAC-learnable but not efficiently exactly learnable) also holds, because simulating an equivalence oracle in the PAC setting requires enumerating hypotheses and checking consistency, which may be computationally intractable for some classes.

\begin{remark}[The role of the cryptographic assumption]
\label{rmk:crypto-assumption}
The Kearns--Valiant result is \emph{conditional}: it assumes that certain cryptographic primitives are hard to break. Without this assumption, the PAC-hardness of DFAs is open. This is typical of computational hardness results in learning theory: unconditional lower bounds for PAC learning remain rare, and most known hardness results reduce from cryptography or from $\mathsf{NP}$-hardness assumptions. The exact learning model sidesteps computational hardness entirely; query access changes the information-theoretic landscape so fundamentally that the cryptographic barrier does not apply.
\end{remark}

\begin{figure}[ht]
\centering
\begin{tikzpicture}[
    box/.style={draw, thick, rounded corners, minimum width=2.6cm, minimum height=0.9cm, font=\small\bfseries, align=center},
    arr/.style={-{Stealth[length=2mm]}, thick},
    noarr/.style={thick, thmred, dashed},
    note/.style={font=\scriptsize, align=center, text=notegray}
]
\node[box, fill=defblue!10] (exact) {Exact Learning\\{\scriptsize (MQ + EQ)}};
\node[box, fill=edgeorange!10, right=4cm of exact] (pac) {PAC Learning\\{\scriptsize (random examples)}};
\node[box, fill=sepgreen!10, below=2.5cm of $(exact)!0.5!(pac)$] (dfa) {DFAs};

% Arrows
\draw[noarr] (pac) -- node[above, font=\scriptsize, text=thmred] {$\not\Rightarrow$} (exact);
\draw[noarr] (exact) -- node[above, font=\scriptsize, text=thmred] {$\not\Rightarrow$} (pac);

\draw[arr, sepgreen!60!black] (dfa) -- node[left, note, text width=2.8cm] {Learnable:\\$L^*$ in poly queries\\{\tiny (Angluin 1987)}} (exact);
\draw[arr, thmred!60!black, dashed] (dfa) -- node[right, note, text width=2.8cm] {Not learnable:\\crypto hardness\\{\tiny (Kearns--Valiant 1994)}} (pac);
\end{tikzpicture}
\caption{The passive--active separation. DFAs are exactly learnable with query access but not PAC-learnable under cryptographic assumptions. Neither paradigm implies the other.}
\label{fig:passive-active-gap}
\end{figure}

% ================================================================
% SECTION 4: QUERY COMPLEXITY AND FORWARD REFERENCES
% ================================================================

\subsection{Extracting an automaton from a transformer}
\label{sec:dfa-extraction}

The passive--active gap has a direct modern use: it is the principle behind \emph{automaton extraction} from trained neural sequence models.  A trained attention model computes a binary string function (accept/reject, or a thresholded next-token decision), and we may treat that function as the unknown target of an exact-learning game, with the model itself playing the minimally adequate teacher.

\begin{prop}[A trained transformer is an $L^*$ target]
\label{prop:transformer-extraction}
Let $f$ be the binary string function computed by a trained attention model.  Answer membership queries by a forward pass, $\mathrm{MQ}(s) = f(s)$, and equivalence queries by sampling-based testing of a conjectured DFA against $f$.  Running $L^*$ against this teacher (\cref{thm:lstar-complexity}) yields a DFA consistent with $f$ on the queried strings, and the minimal DFA for $f$ if the equivalence oracle is exact.\formal{FLT_Proofs.Computation.DFA'} The teacher is a minimally adequate teacher in the sense of \cref{def:mat}, with the model's forward pass as the membership oracle.\formal{FLT_Proofs.Learner.Active.MinimallyAdequateTeacher}
\end{prop}

\begin{proof}
The membership oracle $\mathrm{MQ}(s) = f(s)$ is total and truthful: a forward pass computes $f(s)$ exactly up to the bounded finite-precision forward error of the attention layer.\formal{TLT.Fp32AttnLit.attnLiteralForwardError} The equivalence oracle is approximate: it certifies $L(M) = f$ on a sampled test set rather than on all strings, so the conjecture is exact on the queries made and probably-approximately correct off them; with an exact equivalence oracle $L^*$ returns the minimal DFA for $f$ by \cref{thm:lstar-complexity}.  Either way the model is a legitimate $L^*$ teacher.
\end{proof}

This is how Weiss, Goldberg, and Yahav~\cite{WeissGoldbergYahav2018} extract interpretable automata from recurrent networks, and the construction applies to attention models verbatim.  It turns the passive--active gap into a methodology: a model opaque under passive inspection becomes a queryable teacher, and $L^*$ distills its regular core into a DFA one can read.  The caveat is the one above: the equivalence oracle is only approximate, so extraction is sound on the queries and heuristic beyond them, and whether a transformer's string function even \emph{has} a finite regular core is itself open (\cref{op:transformer-regular}).

\section{Query Complexity}
\label{sec:query-complexity}

The $L^*$ algorithm's complexity is measured in \emph{query complexity}: the number of membership and equivalence queries as a function of the target's representation size. This measure is the exact learning analogue of sample complexity in PAC learning.

\begin{defn}[Query Complexity]
\label{def:query-complexity}
The \emph{membership query complexity} $\mathrm{MQ}(\mathcal{C}, n)$ and \emph{equivalence query complexity} $\mathrm{EQ}(\mathcal{C}, n)$ of a class $\mathcal{C}$ are the minimum numbers of membership and equivalence queries, respectively, that any exact learning algorithm must make in the worst case to identify any target $c \in \mathcal{C}$ of representation size $n$.
\end{defn}

For DFAs, $L^*$ achieves $\mathrm{EQ} \leq n - 1$ and $\mathrm{MQ} = O(n^2 |\Sigma| m)$. The equivalence query count has a direct connection to the combinatorial structure of the hypothesis class.

\begin{prop}[Equivalence queries alone are expensive {\cite{Angluin1988}}]
\label{prop:eq-lower}
With equivalence queries alone (no membership queries), exactly learning a finite class of $N$ concepts requires at most $N - 1$ queries always, and $N - 1$ in the worst case. Membership queries break this bound: $L^*$ uses only $n - 1$ equivalence queries for an $n$-state DFA, although $|\mathcal{C}_n|$ is exponential in $n$.
\end{prop}

\begin{proof}
Maintain the version space $V$ of candidates consistent with every counterexample received, initially the whole class.  Each query $\mathrm{EQ}(h)$ either returns \textsc{yes} (and the learner halts) or a counterexample that is inconsistent with $h$, so it removes at least $h$ from $V$; hence $|V|$ drops by at least one per query and $N - 1$ queries suffice to reach $|V| = 1$.  For the matching lower bound, take a class in which each concept differs from every other at a point unique to that pair, so that for any conjecture $h$ the adversary can return a counterexample consistent with all surviving candidates except $h$; then exactly one candidate is eliminated per query and $N - 1$ are forced.  The contrast with $L^*$ is the content of the chapter: membership queries supply the structural information (which Myhill--Nerode class a string belongs to) that equivalence queries alone cannot, collapsing an exponential query count to a linear one (\cref{prop:eq-alone} of \Cref{ch:data}).
\end{proof}

\begin{remark}[Certificate complexity]
\label{rmk:certificate}
The exact learning framework connects to \emph{certificate complexity} in Boolean function theory: the minimum number of input bits that must be queried to certify a function's value. Teaching dimension (\Cref{ch:dimensions}), which measures the minimum number of examples a teacher must provide to uniquely identify a concept, provides a lower bound on the number of equivalence queries. The connections among query complexity, teaching dimension, and certificate complexity are explored in \Cref{ch:dimensions,ch:computational}.
\end{remark}

\section{What This Chapter Established}
\label{sec:ch8-summary}

\begin{enumerate}[leftmargin=*, itemsep=3pt]
\item \textbf{The exact learning framework.} A minimally adequate teacher provides membership and equivalence queries. Exact learning requires identifying the target concept precisely, not approximately.

\item \textbf{The $L^*$ algorithm.} Angluin's algorithm learns the minimal DFA for any regular language using polynomially many queries. The observation table approximates the Myhill--Nerode equivalence relation; each counterexample refines it by introducing a new state distinction.

\item \textbf{The passive--active gap.} DFAs are exactly learnable with MQ~+~EQ (Angluin, 1987) but not PAC-learnable under cryptographic assumptions (Kearns--Valiant, 1994). This separation is qualitative, not quantitative: query access changes the complexity landscape fundamentally.

\item \textbf{Query complexity.} The number of queries required to exactly learn a class is the natural complexity measure, analogous to sample complexity in PAC learning.
\end{enumerate}

The gap between passive and active learning demonstrated by DFAs is not an isolated curiosity. It reflects a structural principle: the information geometry of a learning problem can change qualitatively when the learner moves from passive observation to active interrogation. This principle recurs in statistical experimental design, active learning in the PAC setting, and reinforcement learning: all settings where the learner's ability to choose which data to collect transforms the computational landscape.

\section{Open Problems}
\label{sec:ch8-open}

Each problem is anchored to a result above and tagged: a \emph{known unknown} (\textsc{ku}) is an articulated open question; an \emph{unknown known} (\textsc{uk}) is a result in the literature or the verified development not yet absorbed into a clean statement.  Verified handles are named where they exist.

\begin{openprob}[The regular core of a transformer, \textsc{uk}]
\label{op:transformer-regular}
\Cref{prop:transformer-extraction} extracts a DFA from a trained attention model, exact only on the queried strings.  Determine when a transformer's binary string function has a finite Myhill--Nerode index (a genuine regular core) so that $L^*$ extraction is exact rather than heuristic, and bound the number of states as a function of the architecture (depth, heads, context length).
\end{openprob}

\begin{openprob}[An exact equivalence oracle for neural models, \textsc{ku}]
\label{op:exact-eq-oracle}
Extraction is approximate because the equivalence oracle samples rather than verifies.  Determine for which model classes an \emph{exact} equivalence query (``does this DFA agree with the model on all strings?'') is decidable, turning approximate extraction into exact $L^*$ learning, and relate the cost to formal verification of the attention layer.
\end{openprob}

\begin{openprob}[Formalizing $L^*$, \textsc{uk}]
\label{op:lstar-formal}
The $L^*$ correctness and complexity theorem (\cref{thm:lstar-complexity}) is cited, not verified.  Formalize $L^*$ (the observation table, closedness, consistency, and the Myhill--Nerode termination argument) on top of the typed minimally adequate teacher (\leanname{FLT_Proofs.Learner.Active.MinimallyAdequateTeacher}) and \leanname{FLT_Proofs.Computation.DFA'}, with the $n-1$ equivalence-query bound proved.
\end{openprob}

\begin{openprob}[A combinatorial dimension for MAT-learnability, \textsc{ku}]
\label{op:mat-dimension}
Which classes are exactly learnable from a minimally adequate teacher in polynomial time?  Beyond DFAs (decision trees, certain Boolean formulas, residual automata are known), give a combinatorial dimension that characterizes polynomial-query MAT-learnability the way $\VC$ characterizes PAC learnability, or prove the query model admits none.
\end{openprob}

\begin{openprob}[The reverse separation, \textsc{ku}]
\label{op:reverse-separation}
The passive--active gap (\cref{sec:passive-active-gap}) is verified in one direction (DFAs: exact yes, PAC no).  Make the reverse precise: exhibit a class that is efficiently PAC-learnable but not efficiently exactly learnable, pinning down why simulating an equivalence oracle in the PAC setting is computationally intractable for it.
\end{openprob}

\begin{openprob}[Sharp equivalence-query complexity, \textsc{ku}]
\label{op:eq-sharp}
\Cref{prop:eq-lower} gives the worst-case $N-1$ bound for equivalence queries alone.  Characterize the exact equivalence-query complexity of a class by a combinatorial parameter (a ``certificate'' or ``extended teaching'' dimension), interpolating between the $N-1$ of unstructured classes and the $n-1$ of DFAs under added membership queries.
\end{openprob}

\begin{openprob}[Query complexity versus teaching dimension, \textsc{uk}]
\label{op:query-teaching}
\Cref{rmk:certificate} connects equivalence queries to teaching dimension and certificate complexity.  Make the connection quantitative: prove two-sided bounds relating $\mathrm{EQ}(\mathcal{C}, n)$ to the teaching dimension (\leanname{FLT_Proofs.Learner.Active.IsOptimalTeacher}) and the certificate complexity of $\mathcal{C}$, unifying the active-learning and teaching pictures of \Cref{ch:dimensions}.
\end{openprob}

\begin{openprob}[Test complexity of attention extraction, \textsc{uk}]
\label{op:extraction-samples}
The approximate equivalence oracle of \cref{prop:transformer-extraction} samples strings to test a conjecture.  Determine the number of test samples needed to certify $L(M) = f$ with confidence $1 - \delta$ as a function of the DFA size and the model's margin, giving a PAC guarantee on the extracted automaton.
\end{openprob}

\begin{openprob}[Counterexample length, \textsc{ku}]
\label{op:counterexample-length}
The $L^*$ membership-query count carries a factor $m$, the longest counterexample length (\cref{rmk:counterexample-length}).  Determine whether a learner or a cooperative teacher can avoid the linear dependence on $m$, for instance by binary-searching counterexamples (Rivest--Schapire), and what the optimal dependence is.
\end{openprob}

\begin{openprob}[An LLM as a minimally adequate teacher, \textsc{uk}]
\label{op:llm-teacher}
The development types an LLM critic as a teacher (\leanname{FLT_Proofs.Learner.Active.LLMCritic}).  Determine the soundness conditions under which a large language model can serve as the membership and equivalence oracle for a downstream exact learner (when its answers are reliable enough that $L^*$ converges to the intended concept), bounding the error introduced by an unreliable teacher.
\end{openprob}


\subsection*{Counterfactual exercises: generalizing Figure~\ref{fig:lstar-protocol}}

Each exercise perturbs the protocol. Show, in the figure's grammar (a cyclic state-discovery loop in which a learner refines a finite approximation of the target's Myhill--Nerode right-congruence by membership queries and is corrected by counterexamples from an equivalence oracle, with the flowchart's nodes drawn solid where a verified data-type certifies the object and the correctness-and-termination claims drawn dashed where they are Angluin's theorem, cited not proved), that the generalized picture subsumes the concrete one as an \emph{instance} (the loop run on a determined target), a \emph{boundary} (where one capability of the teacher or one condition on the table is toggled and the learnable class strictly grows or shrinks), or a \emph{frontier} (a dashed, open or literature-only edge).

\begin{enumerate}[leftmargin=*,itemsep=3pt]
\item \textbf{The table is the right congruence.} Read the observation table not as scratch space but as the object being learned. Show this is the figure's defining instance: two access strings share a row exactly when they are indistinguishable by the suffixes in $E$, a finite slice of the Myhill--Nerode relation $s_1 \sim s_2 \iff (\forall w,\ s_1 w \in L \Leftrightarrow s_2 w \in L)$ (\cref{rmk:table-intuition}), so the distinct rows \emph{are} the equivalence classes discovered so far and the conjecture DFA is their quotient; the row map is typed by the truthful membership oracle (\leanname{FLT_Proofs.Data.MembershipOracle}).

\item \textbf{Closedness and consistency are the quotient's well-definedness.} Toggle off closedness: let some $\mathrm{row}(s\cdot a)$ match no representative in $S$. Show this is the boundary that makes the conjecture illegal: without closedness the transition $\delta(\mathrm{row}(s),a)=\mathrm{row}(s\cdot a)$ has no target state, and without consistency it depends on the chosen representative (\cref{def:closed-consistent}); the two conditions are precisely what licenses the quotient automaton (\leanname{FLT_Proofs.Computation.DFA'}), so they are well-definedness constraints, not optimizations.

\item \textbf{A counterexample is a new-class witness.} Replace the teacher's \textsc{yes}/counterexample reply by a bare bit ``correct or not''. Show this is the boundary that disables refinement: the equivalence oracle's counterexample carries the contract that $z$ is a point where the conjecture and target disagree (\leanname{FLT_Proofs.Data.EquivalenceOracle}; the witness is typed by \leanname{FLT_Proofs.Data.Counterexample}), and adding $z$ with its prefixes forces at least one new distinct row, i.e.\ splits a Myhill--Nerode class; a bare bit removes the witness and the loop can no longer discover the missing state.

\item \textbf{The $n-1$ bound is the index of the congruence.} Read ``at most $n-1$ equivalence queries'' not as a runtime figure but as a structural identity. Show this is the frontier that grades termination: the minimal DFA has $n$ states, hence the congruence has index $n$ (\cref{rmk:table-intuition}, uniqueness by \cite{RabinScott1959}), the table starts at one class and gains at least one per counterexample, so termination after $\le n-1$ rounds \emph{is} the finiteness of the index (\cref{thm:lstar-complexity}); the bound is a literature theorem (\cite{Angluin1987}), not yet kernel-proved (\cref{op:lstar-formal}), so the termination annotation is drawn dashed.

\item \textbf{The whole loop is cited, not verified.} Ask the kernel for the statement ``$L^*$ converges to the minimal DFA''. Show this is the figure's principal frontier: the nodes are verified types ($L^*$ itself is the structure \leanname{LStar} with fields \texttt{teacher}, \texttt{numStates}, \texttt{result}, all sorry-free, but with no method and no correctness lemma), so the objects are solid while the convergence-and-termination edges are dashed and carry \cite{Angluin1987}; closing this gap on top of the typed teacher (\leanname{FLT_Proofs.Learner.Active.MinimallyAdequateTeacher}) and \leanname{FLT_Proofs.Computation.DFA'} is exactly \cref{op:lstar-formal}.

\item \textbf{The teacher is the only source of state-splits.} Read the minimally adequate teacher as the entire interface between learner and target. Show this is the instance that names the model: a MAT is a membership oracle and an equivalence oracle answering about one shared target (\leanname{FLT_Proofs.Learner.Active.MinimallyAdequateTeacher}, field \texttt{consistent} forcing $\texttt{mq.target}=\texttt{eq.target}$; \cref{def:mat}), and nothing else, so every state the learner discovers is forced by a query, never by distributional luck; this is what makes exact identification (\leanname{FLT_Proofs.Criterion.PAC.ExactLearnable}, \cref{def:exact-learning-formal}) attainable here.

\item \textbf{Drop membership queries and the loop blows up.} Run the same cycle with equivalence queries alone. Show this is the boundary that collapses the discovery rate: without membership queries the learner cannot place a string in its Myhill--Nerode class, so each equivalence query removes only one candidate from the version space and an $N$-concept class costs $N-1$ queries in the worst case (\cref{prop:eq-lower}, \cite{Angluin1988}), whereas membership queries supply the structural bit that collapses an exponential count to the linear $n-1$ of \cref{thm:lstar-complexity}.

\item \textbf{Passive data cannot drive the loop at all.} Replace the active teacher with random labelled examples drawn from a distribution. Show this is the frontier that separates the Wills: the same target class (DFAs) is exactly learnable by this loop (\cref{thm:lstar-complexity}) yet not PAC-learnable under cryptographic assumptions (\cref{thm:dfa-pac-hard}, \cite{KearnsValiant1994}), so the passive--active edge is qualitative, and the reverse direction (PAC-easy yet exact-hard) is itself open (\cref{op:reverse-separation}).

\item \textbf{The counterexample length is a hidden cost on the arrows.} Refine the $z$-arrow by asking how long $z$ may be. Show this is a frontier on the membership budget: the equivalence oracle may return an adversarially long counterexample (\cref{rmk:counterexample-length}), so the membership-query count carries a factor $m$ (\cref{thm:lstar-complexity}); whether a learner can binary-search the counterexample to remove the linear dependence on $m$ is open (\cref{op:counterexample-length}, \cite{RivestSchapire1993}).

\item \textbf{Any queryable function is an $L^*$ target.} Replace the unknown regular language by the binary string function of a trained attention model, answering membership by a forward pass and equivalence by sampling. Show this is the figure's most general instance: the model is a legitimate minimally adequate teacher (\cref{prop:transformer-extraction}, \cite{WeissGoldbergYahav2018}), so the same state-discovery loop extracts a DFA exact on the queried strings; whether the model's function even has a finite Myhill--Nerode index (a genuine regular core) so that extraction is exact rather than heuristic is open (\cref{op:transformer-regular}), and so is the soundness of an unreliable oracle such as an LLM critic (\cref{op:llm-teacher}).
\end{enumerate}
